\begin{lemma}[A Tashkinov tree is generated from its one-edge prefix]
Let \(G\) be a finite loopless multigraph with a finite edge-identity set,
let \(e_0\) be an edge identity, let \(\varphi\) be a partial \(k\)-edge-coloring,
and let \(T=(y_0,e_1,y_1,\ldots,e_n,y_n)\) be a Tashkinov tree with
respect to \(e_0\) and \(\varphi\).  Then there is a Tashkinov augmentation
sequence under \(\varphi\) from the exact length-one prefix
\(T|y_1=(y_0,e_1,y_1)\) to \(T\).
\end{lemma}
\begin{proof}
For every \(q\) with \(1\leq q\leq n\), let \(P_q=T|y_q\), using the
exact prefix from \noderef{TreeSequencePrefix}.  Thus \(P_q\) has length
\(q\), its displayed vertices are exactly \(y_0,\ldots,y_q\) in that
order, and its displayed edge identities are exactly
\(e_1,\ldots,e_q\) in that order.

Fix \(q\) with \(1\leq q<n\).  Apply the defining property in
\noderef{TashkinovTree} to the next displayed edge \(e_{q+1}\).  It gives
a color \(c\) such that \(\varphi(e_{q+1})=c\) and an index
\(i<q+1\) such that \(c\) is missing at \(y_i\).  Since
\(i\leq q\), the vertex \(y_i\) belongs to \(P_q\), so \(c\) belongs to
the union of the colors missing at vertices of \(P_q\), namely the tree
missing-color set from \noderef{TreeMissingColors}.

The prefix \(P_{q+1}\) has length \(q+1\), contains \(P_q\) as its exact
ordered initial segment, and its newly added last edge is the same edge
identity \(e_{q+1}\), not merely an edge with the same endpoints.  That
edge has color \(c\), which is missing on \(P_q\).  These are precisely
the clauses of \noderef{TashkinovAugmentation}; hence \(P_{q+1}\) is one
Tashkinov augmentation of \(P_q\) under \(\varphi\).

Starting with \(P_1\), induction on \(q\) and concatenation of the preceding
one-step augmentations gives a Tashkinov augmentation sequence from \(P_1\)
to every \(P_q\).  The base case is the empty augmentation sequence, allowed
by the reflexive clause in \noderef{TashkinovAugmentationSequence}.  At
\(q=n\), the exact full-length prefix has all the vertex and edge data of
\(T\), hence is \(T\) itself.  The resulting sequence from \(P_1\) to \(T\)
is the required one.
\end{proof}
